\textit{Catatan edisi: solusi sumber memakai reduksi ke matriks Hesse tanpa menyebut matriks metrik grafik dan memberikan vektor eigen nol pada kasus sah $c=0$ dan $a<b$. Alasan metrik dan pemisahan kasus berikut memperbaiki kedua kekurangan itu.}

Kita menggunakan
[[Differenzierbare Funktion/Graph/Hyperfläche/Weingartenabbildung/Fakt|Fakta]],
dan di titik nol berlaku $\operatorname{Grad}f(0)=0$, sehingga faktor normalisasi $\omega=1$ dan matriks metrik grafik $G=I$. Oleh karena itu,
\definitionsverweis {pemetaan Weingarten}{}{}
langsung dinyatakan oleh
\definitionsverweis {matriks Hesse}{}{.}
Matriks Hesse dari fungsi tersebut adalah

\mathdisp {\begin{pmatrix}  2a  &   c   \\ c &  2b \end{pmatrix}} { , }

sedangkan
\definitionsverweis {polinom karakteristik}{}{}
darinya adalah

\mavergleichskettealign
{\vergleichskettealign
{ (t-2a)(t-2b)-c^2
}
{ =} { t^2 -2(a+b)t +4ab -c^2
}
{ =} { \left(  t-(a+b)  \right)^2 - (a+b)^2 +4ab -c^2
}
{ =} { \left(  t-(a+b)  \right)^2  - a^2 -b^2 +2ab -c^2
}
{ =} { \left(  t-(a+b)  \right)^2  - (a-b)^2 -c^2
}
}
{}
{}{.}
Jadi, nilai-nilai eigennya
\zusatzklammer {yang sekaligus merupakan kelengkungan utama} {} {}
adalah

\mavergleichskettedisp
{\vergleichskette
{ \lambda_{1,2}
}
{ =} { \pm \sqrt{ (a-b)^2+c^2 } +a+b
}
{ } {
}
{ } {
}
{ } {
}
}
{}{}{.}
Jika $c=0$, matriks tersebut diagonal. Apabila $a>b$, arah eigen
\mathkor {} {\begin{pmatrix}1\\0\end{pmatrix}} {dan} {\begin{pmatrix}0\\1\end{pmatrix}} {}
masing-masing bersesuaian dengan nilai eigen $2a$ dan $2b$. Apabila $a<b$, arah eigen
\mathkor {} {\begin{pmatrix}0\\1\end{pmatrix}} {dan} {\begin{pmatrix}1\\0\end{pmatrix}} {}
masing-masing bersesuaian dengan nilai eigen $2b$ dan $2a$. Jika

\mavergleichskette
{\vergleichskette
{ a
}
{ = }{ b
}
{  }{
}
{    }{
}
{   }{
}
}
{}{}{}
dan

\mavergleichskette
{\vergleichskette
{ c
}
{ = }{ 0
}
{  }{
}
{    }{
}
{   }{
}
}
{}{}{}
maka pemetaan Weingarten adalah perkalian dengan $2a$ dan setiap vektor tak nol merupakan vektor eigen. Dengan demikian, seluruh kasus $c=0$ telah selesai. Mulai sekarang, anggap $c\neq 0$. Untuk

\mavergleichskettedisp
{\vergleichskette
{ \lambda_{1}
}
{ =} { \sqrt{ (a-b)^2+c^2 } +a+b
}
{ } {
}
{ } {
}
{ } {
}
}
{}{}{}
kita perlu menentukan kernel matriks

\mathdisp {\begin{pmatrix}  \sqrt{(a-b)^2 +c^2} + b-a  &   -c  \\  -c &  \sqrt{(a-b)^2 +c^2} + a-b \end{pmatrix}} {  }

yang diberikan oleh

\mathdisp {\R \begin{pmatrix}  \sqrt{(a-b)^2 +c^2} + a-b  \\c   \end{pmatrix}} { . }

Dengan demikian, salah satu arah kelengkungan utama adalah
\mathl{\begin{pmatrix}  \sqrt{(a-b)^2 +c^2} + a-b  \\c   \end{pmatrix}}{} dengan kelengkungan utama
\mathl{\sqrt{ (a-b)^2+c^2 } +a+b}{.}

Dengan cara yang sama, diperoleh arah kelengkungan utama
\mathl{\begin{pmatrix}  -c \\ \sqrt{(a-b)^2 +c^2} + a-b    \end{pmatrix}}{} dengan kelengkungan utama
\mathl{-\sqrt{ (a-b)^2+c^2 } +a+b}{.}

 [[Kategorie:Latexseite]]
